\documentclass[10pt,a4paper]{amsart}
\usepackage[margin=19mm]{geometry}
\usepackage{amsmath,amssymb,mathtools}
\usepackage[scr=rsfs,cal=euler]{mathalfa}
\usepackage{xfrac}
\usepackage[unicode,hidelinks,bookmarks=false]{hyperref}
\newcommand{\ie}{i.e.\ }
\newcommand{\cf}{cf.\ }
\newcommand{\resp}{respectively\ }
\newcommand{\Subsection}{\subsection}
\providecommand{\nobreakdash}{\nobreak\hbox{-}}
\setlength{\emergencystretch}{2em}
\multlinegap0pt
\usepackage{bm}
\usepackage{bbm}
\usepackage{dsfont}
\usepackage{xspace}
\usepackage{cancel}
\usepackage{mathtools}
\usepackage{amsthm}
\makeatletter
\@ifundefined{Theorem}{\newtheorem{Theorem}{Theorem}}{}
\makeatother
\newcommand{\card}[1]{\left|#1\right|}
\newcommand{\mcl}[1]{\ensuremath{\mathcal{#1}}}
\title[The Kirkwood-Dirac representation associated to the Fourier ]{The Kirkwood-Dirac representation associated to the Fourier transform for finite abelian groups: positivity}
\author{Stephan De Bi\`evre, Christopher Langrenez, Danylo Radchenko}
\date{Source: arXiv:2501.12252v2 (CC BY 4.0)}
\begin{document}
\maketitle

\noindent\emph{Source and licence.} The Kirkwood-Dirac representation associated to the Fourier transform for finite abelian groups: positivity. Version arXiv:2501.12252v2 (CC BY 4.0), distributed under the Creative Commons Attribution 4.0 International licence, as recorded in the arXiv licence metadata for this exact version and shown on the abstract page https://arxiv.org/abs/2501.12252v2. Full rights notice: licenses/arxiv-2501.12252-CC-BY-4.0.txt.\par\smallskip

\noindent\textit{Editorial scope.} Counted unit: the Theorem whose source label is \texttt{thm:KDPure} in the article (source0.tex lines 446--448, source label \texttt{thm:KDPure}), delivered together with the article's own complete original proof of it (source0.tex lines 450--494, one \texttt{proof} environment of 45 lines). No mathematics is written, repaired or re-derived: statement and proof are the article's own text line for line. Required context retained uncounted: none. Every definition, display and label of those lines is retained unchanged. Omitted from this package and not redistributed: 1--445, 449--449, 495--790. Nothing inside a retained excerpt is omitted or altered: every retained body is delivered exactly as the article's lines are written, so no omission receipt of any kind accompanies this package. Citations: the retained excerpts carry no citation command at all, so no bibliography entry of the article is transcribed and no reference is redistributed. Reference adaptation: none was needed and none was made; every reference inside the delivered unit resolves inside it, so no unresolved reference or citation is produced. Printed slips kept literal: no printed word, symbol, hypothesis or number of the counted statement or of its proof was altered. Layout and class: the article's own class is \texttt{article} with packages including inputenc, babel, fontenc, amsmath, amsthm, amssymb, amscd, physics, color, savesym, soul, url, ulem, enumitem; the wrapper is the lane's pinned \texttt{amsart} editorial wrapper at 10pt on A4, which restates the theorem-like environments the retained text needs and adds the article's own macro definitions that the retained text needs (\texttt{\textbackslash card}, \texttt{\textbackslash mcl}), with extra wrapper packages (bm, bbm, dsfont, xspace, cancel, mathtools). The delivered numbering is the unit's own. Assets: no retained span contains a figure environment, a graphics-inclusion command, a PDF-page inclusion, a table or a TikZ picture; no image file is copied into this package, the package declares no asset dependency and no image file is redistributed with it. Licence: version arXiv:2501.12252v2 of this article is distributed under the Creative Commons Attribution 4.0 International licence, as recorded in the arXiv licence metadata for this exact version and shown on the abstract page https://arxiv.org/abs/2501.12252v2. A full rights notice is delivered at licenses/arxiv-2501.12252-CC-BY-4.0.txt. Characters: every retained excerpt is pure ASCII, so the delivered engine, XeLaTeX with its default Latin Modern text font, reports no missing character.
\begin{Theorem}\label{thm:KDPure}
     Let $G$ be a finite abelian group. Then $\psi$ is a pure KD-positive state if and only if there exists a subgroup $H$ and $(g_0,\chi_0)\in G/H \times\hat{G}/H^{\perp}$ such that $\psi = \psi^{H}_{g_0,\chi_0}$.
 \end{Theorem}

\begin{proof}
    As $\psi^{H}_{g_0,\chi_0}$ is a KD-positive pure state for all subgroups $H$ and all $(g_0,\chi_0)\in G/H \times\hat{G}/H^{\perp}$, we only need to prove the direct implication. 
    
    Let $\psi$ be a KD-positive state. We denote by 
    \begin{equation}
        S=\mathrm{supp}(\psi):= \left\{g\in G, \psi(g) \neq 0\right\} \mathrm{\ and \ } S' = \mathrm{supp}(\hat{\psi}):= \left\{\chi\in \hat{G}, \hat{\psi}(\chi) \neq 0\right\}.
    \end{equation}
    Thus, for all $g\in S, \chi \in S'$, as $\psi$ is KD-positive, we have that:
    \begin{equation}\label{eq:KDpurepos1}
        \overline{\chi(g)}\psi(g)\overline{\hat{\psi}(\chi)} >0.
    \end{equation}
    We obtain that, for all $(g_1,g_2)\in S^2, \chi \in S'$:
    \begin{equation}\label{eq:KDpurepos2}
        \left(\overline{\chi(g_1)}\psi(g_1)\overline{\hat{\psi}(\chi)}\right)^{-1} >0 \mathrm{ \ and \ } \overline{\chi(g_2)}\psi(g_2)\overline{\hat{\psi}(\chi)} > 0.
    \end{equation}
    Thus, by multiplying the two inequalities in Eq.\eqref{eq:KDpurepos2}, for all $(g_1,g_2)\in S^2, \chi \in S'$
    \begin{equation}\label{eq:KDpurepos3}
        \chi(g_1-g_2)\psi(g_1)^{-1}\psi(g_2) > 0.
    \end{equation}
    Furthermore, by using Eq.\eqref{eq:KDpurepos3}, we obtain that, for all $(g_1,g_2)\in S^2, (\chi_1,\chi_2) \in S'^{2}$:
    \begin{equation}\label{eq:KDpurepos4}
        \chi_{1}(g_1-g_2)\overline{\chi_{2}(g_1-g_2)} > 0.
    \end{equation}
    As the action of the Heisenberg group preserves KD-positivity, we can suppose that $0\in S$ and $1\in S'$. We apply Eq.\eqref{eq:KDpurepos4} with $g_2=0$ and $\chi_2=1$ to obtain:
    \begin{equation}\label{eq:KDpurepos5}
        \forall (g,\chi) \in S\times S', \chi(g) > 0.
    \end{equation}
    As for all $(g,\chi) \in S\times S', \chi(g)\in\mcl{S}^{1}$, we have, with Eq.\eqref{eq:KDpurepos5}, that $ \chi(g)=1$. This implies that $S'\subseteq S^{\perp}$. Consequently, as $\left|S^{\perp}\right| \leqslant  \frac{\left|G\right|}{\left|S\right|}$, we obtain that 
    \begin{equation}\label{eq:KDpurepos6}
        \card{S}\card{S'} \leqslant \card{G}.
    \end{equation}
    Moreover, we have that:
    \begin{equation}
             \card{G}=\card{G}\sum_{(g,\chi)\in G\times\hat{G}} \left|Q[\psi](g,\chi)\right| = \card{G}^{\frac{1}{2}}\sum_{(g,\chi)\in G\times\hat{G}} \left|\psi(g)\hat{\psi}(\chi)\right| = \card{G}^{\frac{1}{2}}\sum_{g\in S} \left|\psi(g)\right| \sum_{\chi\in S'}\left|\hat{\psi}(\chi)\right|.
    \end{equation}
    We apply the Cauchy-Schwarz inequality twice to obtain that
    \begin{equation}
        \card{G}\leqslant \card{G}^{\frac{1}{2}}\left(\sum_{g\in S} \left|\psi(g)\right|^2\right)^{\frac{1}{2}}\card{S}^{\frac{1}{2}}\left(\sum_{\chi\in S'}\left|\hat{\psi}(\chi)\right|^2\right)^{\frac{1}{2}}\card{S'}^{\frac{1}{2}} = \card{G}^{\frac{1}{2}}\card{S}^{\frac{1}{2}}\card{S'}^{\frac{1}{2}} 
    \end{equation}
    Squaring this inequality and reorganising the terms, we have that:
    \begin{equation}\label{eq:KDpurepos7}
        \card{G} \leqslant \card{S}\card{S'}.
    \end{equation}
    Eq.\eqref{eq:KDpurepos6} and Eq.\eqref{eq:KDpurepos7} together imply that $\card{G} = \card{S}\card{S'}$ and thus, $\card{S'} = \card{S^{\perp}}$. Thus $S'=S^{\perp}$ and $S'$ is a subgroup of $\hat{G}$. Similarly, $S$ is a subgroup of $G$. And by equality case in the Cauchy-Schwarz inequality, we obtain that $\psi = \psi^{S}$, meaning that $\psi$ is in the Heisenberg orbit of a normalized characteristic function of a subgroup.  
\end{proof}

\end{document}
